\documentclass[10pt,a4paper]{amsart}
\usepackage[margin=19mm]{geometry}
\usepackage{amsmath,amssymb,mathtools}
\usepackage[scr=rsfs,cal=euler]{mathalfa}
\usepackage{xfrac}
\usepackage[unicode,hidelinks,bookmarks=false]{hyperref}
\newcommand{\ie}{i.e.\ }
\newcommand{\cf}{cf.\ }
\newcommand{\resp}{respectively\ }
\newcommand{\Subsection}{\subsection}
\providecommand{\nobreakdash}{\nobreak\hbox{-}}
\setlength{\emergencystretch}{2em}
\multlinegap0pt
\usepackage{tikz}
\usetikzlibrary{cd}
\usepackage{mathtools}
\usepackage{amssymb}
\usepackage{tikz}
\usepackage{dsfont}
\usepackage[shortlabels]{enumitem}
\newtheorem{assumption}{Assumption}
\newtheorem{proposition}{Proposition}
\newtheorem{lemma}{Lemma}
\def\FF{\mathbf{F}}
\def\GL{\mathrm{GL}}
\def\QQ{\mathbf{Q}}
\def\Qbar{\QQ^{\ac}}
\def\ac{\mathrm{ac}}
\def\calO{\mathcal{O}}
\def\coker{\mathrm{coker}}
\def\frakm{\mathfrak{m}}
\def\frakp{\mathfrak{p}}
\def\ram{\mathrm{ram}}
\def\rmG{\mathrm{G}}
\def\rmH{\mathrm{H}}
\def\rmR{\mathrm{R}}
\def\rmT{\mathrm{T}}
\def\rmX{\mathrm{X}}
\def\rmZ{\mathrm{Z}}
\def\ss{\mathrm{ss}}
\title{Arithmetic level raising on triple product of Shimura curves and Gross--Kudla--Schoen Diagonal cycles II: Bipartite Euler system}
\author{Haining Wang}
\date{Source: arXiv:2004.14916v4 (CC BY 4.0)}
\begin{document}
\maketitle

\noindent\emph{Source and licence.} Arithmetic level raising on triple product of Shimura curves and Gross--Kudla--Schoen Diagonal cycles II: Bipartite Euler system. Version arXiv:2004.14916v4 (CC BY 4.0), distributed by arXiv under the Creative Commons Attribution 4.0 International licence, http://creativecommons.org/licenses/by/4.0/, as recorded in the arXiv licence metadata for this exact version and shown on the abstract page https://arxiv.org/abs/2004.14916v4. Full licence notice: \texttt{licenses/arxiv-2004.14916-CC-BY-4.0.txt}.\par\smallskip

\noindent\textit{Editorial scope.} Counted unit: the article's proposition new-old-part (source0.tex lines 1072--1083). The article's complete original proof of it is delivered with it (source0.tex lines 1085--1171, one \texttt{proof} environment of 87 lines). No mathematics is written, repaired or re-derived: the statement and the proof are the article's own lines, in the article's own notation, and no retained line is substituted or re-flowed.  Retained uncounted as the source-local context the proof needs: lines 840--847; lines 899--901; lines 1043--1049.  Nothing was omitted from any retained span, so the package carries no omission record.  No retained line contains a citation, so the wrapper carries no bibliography and no bibliography entry.  The retained lines contain the article's own diagram code, which the wrapper typesets with the standard tikz packages; no image file is delivered and the package declares no asset dependency.  Editorial formatting only, in addition: an AMS \texttt{amsart} preamble that restates the article's own declarations and macros used by the retained lines, namely the environments \texttt{assumption}, \texttt{proposition}, \texttt{lemma} on separate counters, together with the standard packages those lines need.  The retained environments print in this document as Assumption 1, Proposition 1, Lemma 1 in order of appearance, whereas the article prints the same items with the numbering of its own full document.  The complete native source of the article as distributed in the arXiv e-print for this exact version is bundled unchanged as \texttt{sources/103642/source0.tex}.
\begin{assumption}\label{assump}
We make the following assumptions on $\overline{\rho}_{f, \lambda}$ and thus on the maximal ideal $\frakm$:
\begin{enumerate}
\item $\bar{\rho}_{f, \lambda}\vert_{\rmG_{\QQ(\zeta_{\ell})}}$ is absolutely irreducible;
\item $\bar{\rho}_{f, \lambda}$ is minimally ramified at primes in $\Sigma^{+}\cup\Sigma^{-}_{\ram}$ and is ramified at primes in $\Sigma^{-}_{\ram}$;
\item The image of $\overline{\rho}_{f, \lambda}$ contains $\GL_{2}(\FF_{\ell})$.
\end{enumerate}
\end{assumption}

\begin{proposition}\label{Phi-n}
The map $\Phi_{n}$ is surjective and $\Psi_{n}$ is injective.
\end{proposition}

\begin{lemma}\label{canonical-decomposition}
Let $p$ be an $n$-admissible prime for $f$ and suppose that Assumption \ref{assump} holds for $\overline{\rho}_{f,\lambda}$, then we have the following exact sequence of $\calO_{\lambda, n}[\rmG_{\FF_{p^{2}}}]$-modules
\begin{equation*}
0\rightarrow\Gamma(\rmZ_{d}(\overline{B}), \calO_{\lambda}(1))_{/\frakp^{[p]}_{n}}\rightarrow \rmH^{1}(\overline{\rmX}_{d}\otimes{\FF^{\ac}_{p}}, \calO_{\lambda}(1))_{/\frakp_{n}}\rightarrow \Gamma(\rmZ_{d}(\overline{B}), \calO_{\lambda})_{/\frakp^{[p]}_{n}}\rightarrow0
\end{equation*}
which splits.
\end{lemma}

\begin{proposition}\label{new-old-part}
Let $p$ be an $n$-admissible prime for $f$ and suppose that Assumption \ref{assump} holds for $\overline{\rho}_{f,\lambda}$. Then the new part $\rmH^{1}({\rmX}_{d}(p)\otimes{\QQ^{\ac}}, \calO_{\lambda}(1))^{\bullet}_{/\frakp^{[p]}_{n}}$ is unramified at $p$ and we have the following split exact sequence
\begin{equation*}
0\rightarrow\Gamma(\rmZ_{d}(\overline{B}), \calO_{\lambda}(1))_{/\frakp^{[p]}_{n}}\rightarrow \rmH^{1}({\rmX}_{d}(p)\otimes{\QQ^{\ac}_{p}}, \calO_{\lambda}(1))^{\bullet}_{/\frakp^{[p]}_{n}}\rightarrow \Gamma(\rmZ_{d}(\overline{B}), \calO_{\lambda})_{/\frakp^{[p]}_{n}}\rightarrow0
\end{equation*}
of $\calO_{\lambda, n}[\rmG_{\FF_{p^{2}}}]$-modules.
Moreover the map $\alpha^{\ast}=(\pi^{\ast}_{1,p}, \pi^{\ast}_{2,p})$ induces an isomorphism 
\begin{equation*}
\alpha^{\ast}: \rmH^{1}({\rmX}_{d}\otimes{\Qbar}, \calO_{\lambda}(1))_{/\frakp_{n}}\xrightarrow{\sim} \rmH^{1}(\rmX_{d}(p)\otimes{\Qbar}, \calO_{\lambda}(1))^{\bullet}_{/\frakp^{[p]}_{n}}
\end{equation*}
of $\calO_{\lambda, n}[\rmG_{\QQ}]$-modules.
\end{proposition}

\begin{proof}
We have a commutative diagram
\begin{equation}\label{dia1}
\begin{tikzcd}
\ker(\overline{\alpha}_{\ast}) \arrow[d] \arrow[r]            &\ker({\alpha}_{\ast})   \arrow[d] \arrow[r]          &   \rmH^{1}(\overline{\rmX}_{d}(p)\otimes{\FF^{\ac}_{p}}, \rmR\Phi(\calO_{\lambda})(1))_{/\frakp^{[p]}_{n}}    \arrow[d, Rightarrow, no head]       \\
\rmH^{1}(\overline{\rmX}_{d}(p)\otimes{\FF^{\ac}_{p}}, \calO_{\lambda}(1))_{/\frakp^{[p]}_{n}} \arrow[r, "sp"] \arrow[d, "\overline{\alpha}_{\ast}"]      & \rmH^{1}(\overline{\rmX}_{d}(p)\otimes{\FF^{\ac}_{p}}, \rmR\Psi(\calO_{\lambda})(1))_{/\frakp^{[p]}_{n}} \arrow[r] \arrow[d, "\alpha_{\ast}"] & \rmH^{1}(\overline{\rmX}_{d}(p)\otimes{\FF^{\ac}_{p}}, \rmR\Phi(\calO_{\lambda})(1))_{/\frakp^{[p]}_{n}} \arrow[d] \\
\rmH^{1}(\overline{\rmX}_{d}\otimes{\FF^{\ac}_{p}}, \calO_{\lambda}(1))^{\oplus 2}_{/\frakp^{[p]}_{n}}\arrow[r, Rightarrow, no head]\arrow[d] & \rmH^{1}(\overline{\rmX}_{d}\otimes {\FF^{\ac}_{p}}, \calO_{\lambda}(1))^{\oplus 2}_{/\frakp^{[p]}_{n}}\arrow[r] \arrow[d]     & 0      \\
\coker(\overline{\alpha}_{\ast})\arrow[r] &0\\    
\end{tikzcd}
\end{equation}
where $\overline{\alpha}_{\ast}$ and $\alpha_{\ast}$ are both induced by the degeneracy maps $(\pi_{1,p,\ast}, \pi_{2,p,\ast})$. 

To understand $\ker(\overline{\alpha}_{\ast})$, we consider the following commutative diagram
\begin{equation}\label{dia2}
\begin{tikzcd}
\rmH^{0}(\overline{\rmX}^{\ss}_{d}\otimes{\FF^{\ac}_{p}}, \calO_{\lambda}(1))_{/\frakp^{[p]}_{n}}\arrow[r] \arrow[d, Rightarrow, no head] & \ker(\overline{\alpha}_{\ast}) \arrow[d] \arrow[r] & \ker(\nabla_{n}) \arrow[d] \\
\rmH^{0}(\overline{\rmX}^{\ss}_{d}\otimes{\FF^{\ac}_{p}}, \calO_{\lambda}(1))_{/\frakp^{[p]}_{n}} \arrow[r] \arrow[d, "\overline{\alpha}_{\ast}"] & \rmH^{1}(\overline{\rmX}_{d}(p)\otimes{\FF^{\ac}_{p}}, \calO_{\lambda}(1))_{/\frakp^{[p]}_{n}} \arrow[r] \arrow[d, "\overline{\alpha}_{\ast}"] & \rmH^{1}(\overline{\rmX}_{d}\otimes{\FF^{\ac}_{p}}, \calO_{\lambda}(1))^{\oplus 2}_{/\frakp^{[p]}_{n}} \arrow[d, "\nabla_{n}"] \\
0 \arrow[r]  & \rmH^{1}(\overline{\rmX}_{d}\otimes{\FF^{\ac}_{p}}, \calO_{\lambda}(1))^{\oplus 2}_{/\frakp^{[p]}_{n}} \arrow[r, Rightarrow, no head] \arrow[d] & \rmH^{1}(\overline{\rmX}_{d}\otimes{\FF^{\ac}_{p}}, \calO_{\lambda}(1))^{\oplus2}_{/\frakp^{[p]}_{n}} \arrow[d] \\
          & \coker(\overline{\alpha}_{\ast}) \arrow[r]           & \coker(\nabla_{n})     
\end{tikzcd}
\end{equation}
which implies by the snake lemma that $\ker(\overline{\alpha}_{\ast})$ sits in the exact sequence
\begin{equation}\label{ker-alpha}
0\rightarrow \rmH^{0}(\overline{\rmX}^{\ss}_{d}\otimes{\FF^{\ac}_{p}}, \calO_{\lambda}(1))_{/\frakp^{[p]}_{n}}\rightarrow \ker(\overline{\alpha}_{\ast})\rightarrow \ker(\nabla_{n})\rightarrow0
\end{equation}
and that $\coker(\overline{\alpha}_{\ast})=\coker(\nabla_{n})$. By the proof of Proposition \ref{canonical-decomposition}, we have 
\begin{equation*}
\coker(\nabla_{n})\cong\rmH^{1}(\FF_{p^{2}}, \rmH^{1}(\overline{\rmX}_{d}\otimes{\FF^{\ac}_{p}}, \calO_{\lambda}(1))_{/\frakp^{[p]}_{n}}). 
\end{equation*}
Note the connecting map  
\begin{equation*}
\rmH^{1}(\overline{\rmX}_{d}(p)\otimes{\FF^{\ac}_{p}}, \rmR\Phi(\calO_{\lambda})(1))_{/\frakp^{[p]}_{n}}\cong \Gamma(\rmZ_{d}(\overline{B}), \calO_{\lambda})_{/\frakp^{[p]}_{n}}\rightarrow \coker(\overline{\alpha}_{\ast})\cong \rmH^{1}(\FF_{p^{2}}, \rmH^{1}(\overline{\rmX}_{d}\otimes{\FF^{\ac}_{p}}, \calO_{\lambda}(1))_{/\frakp^{[p]}_{n}})
\end{equation*}
is by definition the same as the map $\Phi_{n}$ and hence is an isomorphism.
Therefore by the snake lemma applied to diagram \ref{dia1}, we have 
\begin{equation*}
\ker(\overline{\alpha}_{\ast})=\ker(\alpha_{\ast}).
\end{equation*} 
Note also that $ \ker(\nabla_{n})\cong \coker(\nabla_{n})$ by the canonical isomorphism 
\begin{equation*}
\rmH^{0}(\FF_{p^{2}}, \rmH^{1}(\overline{\rmX}_{d}\otimes{\FF^{\ac}_{p}}, \calO_{\lambda}(1))_{/\frakp^{[p]}_{n}}\cong \rmH^{1}(\FF_{p^{2}}, \rmH^{1}(\overline{\rmX}_{d}\otimes{\FF^{\ac}_{p}}, \calO_{\lambda}(1))_{/\frakp^{[p]}_{n}}
\end{equation*}
which in turn is also isomorphic to $\Gamma(\rmZ_{d}(\overline{B}), \calO_{\lambda}(1))_{/\frakp^{[p]}_{n}}$ via the map $\Phi^{\ast}_{n}$.
Therefore the exact sequence \ref{ker-alpha} is isomorphic to
 \begin{equation*}
0\rightarrow \Gamma(\rmZ_{d}(\overline{B}), \calO_{\lambda}(1))_{/\frakp^{[p]}_{n}}\rightarrow \ker(\overline{\alpha}_{\ast})\rightarrow \Gamma(\rmZ_{d}(\overline{B}), \calO_{\lambda})_{/\frakp^{[p]}_{n}}\rightarrow0
\end{equation*}
which splits as $\ker(\overline{\alpha}_{\ast})$ is clearly unramified as an $\calO_{\lambda, n}[\rmG_{\FF_{p^{2}}}]$-module. This combines with the isomorphism $\ker(\overline{\alpha}_{\ast})=\ker(\alpha_{\ast})$ gives the exact sequence
\begin{equation*}
0\rightarrow \Gamma(\rmZ_{d}(\overline{B}), \calO_{\lambda}(1))_{/\frakp^{[p]}_{n}}\rightarrow \rmH^{1}({\rmX}_{d}(p)\otimes {\QQ^{\ac}_{p}}, \calO_{\lambda}(1))^{\bullet}_{/\frakp^{[p]}_{n}}\rightarrow \Gamma(\rmZ_{d}(\overline{B}), \calO_{\lambda})_{/\frakp^{[p]}_{n}}\rightarrow0.
\end{equation*}

Since the composite map $\alpha_{\ast}\alpha^{\ast}$
\begin{equation*}
\rmH^{1}({\rmX}_{d}\otimes{\Qbar}, \calO_{\lambda}(1))^{\oplus2}_{/\frakp^{[p]}_{n}}\xrightarrow{\alpha^{\ast}} \rmH^{1}(\rmX_{d}(p)\otimes{\Qbar}, \calO_{\lambda}(1))_{/\frakp^{[p]}_{n}}\xrightarrow{\alpha_{\ast}} \rmH^{1}({\rmX}_{d}\otimes{\Qbar}, \calO_{\lambda}(1))^{\oplus2}_{/\frakp^{[p]}_{n}}
\end{equation*}
is given by the matrix 
\begin{equation*}
\Delta_{n}=
\begin{pmatrix}
p+1 &   \rmT_{p}\\
\rmT_{p} & p+1\\ 
\end{pmatrix}
\end{equation*}
whose kernel is isomorphic to $\rmH^{1}({\rmX}_{d}\otimes{\Qbar}, \calO_{\lambda}(1))_{/\frakp^{[p]}_{n}}$ by the map 
\begin{equation}\label{1/2}
x\mapsto \frac{1}{2}(x, -\epsilon_{p}(f)x), 
\end{equation}
we have a canonical map 
\begin{equation*}
\alpha^{\ast}: \ker(\alpha_{\ast}\alpha^{\ast})\cong\rmH^{1}({\rmX}_{d}\otimes{\Qbar}, \calO_{\lambda}(1))_{/\frakp^{[p]}_{n}} \rightarrow \ker(\alpha_{\ast})=\rmH^{1}(\rmX_{d}(p)\otimes{\Qbar}, \calO_{\lambda}(1))^{\bullet}_{/\frakp^{[p]}_{n}} 
\end{equation*}
which is injective by Ihara's lemma. 

Note that Proposition \ref{Phi-n} and Proposition \ref{canonical-decomposition} furnish a map of exact sequences
\begin{equation*}
\begin{tikzcd}
0 \arrow[r] &\Gamma(\rmZ_{d}(\overline{B}), \calO_{\lambda}(1))_{/\frakp^{[p]}_{n}} \arrow[r] \arrow[d]      & \rmH^{1}({\rmX}_{d}\otimes{\QQ^{\ac}_{p}}, \calO_{\lambda}(1))_{/\frakp^{[p]}_{n}} \arrow[r] \arrow[d, "\alpha^{\ast}"] & \Gamma(\rmZ_{d}(\overline{B}), \calO_{\lambda})_{/\frakp^{[p]}_{n}} \arrow[d] \arrow[r] & 0\\
0 \arrow[r] &\Gamma(\rmZ_{d}(\overline{B}), \calO_{\lambda}(1))_{/\frakp^{[p]}_{n}} \arrow[r] &\rmH^{1}({\rmX}_{d}(p)\otimes {\QQ^{\ac}_{p}}, \calO_{\lambda}(1))^{\bullet}_{/\frakp^{[p]}_{n}}\arrow[r] & \Gamma(\rmZ_{d}(\overline{B}), \calO_{\lambda})_{/\frakp^{[p]}_{n}}\arrow[r] & 0\\  
\end{tikzcd}
\end{equation*}
of $\calO_{\lambda,n}[\rmG_{\FF_{p^{2}}}]$-modules. It follows that 
\begin{equation*}
\rmH^{1}(\rmX_{d}(p)\otimes{\Qbar}, \calO_{\lambda}(1))^{\bullet}_{/\frakp^{[p]}_{n}}\cong \rmH^{1}(\rmX_{d}\otimes{\Qbar}, \calO_{\lambda}(1))_{/\frakp_{n}}
\end{equation*}
as they are of the same cardinality by the above discussion and $\alpha^{\ast}$ is injective.
\end{proof}

\end{document}
